\documentclass[12pt]{article}

% Настройка русского языка и шрифтов
\usepackage[utf8]{inputenc}
\usepackage[russian]{babel}
\usepackage[T2A]{fontenc}

% Математические пакеты
\usepackage{amsmath, amssymb, amsfonts}

% Настройка страницы
\usepackage{geometry}
\geometry{a4paper, margin=0.8in}

% Другие полезные пакеты
\usepackage{hyperref}
\usepackage{booktabs}
\usepackage{enumitem}
\usepackage{longtable}
\usepackage{graphicx}
\usepackage{xcolor}

% Определение стилей для комментариев
\definecolor{commentgray}{rgb}{0.4,0.4,0.4}
\definecolor{antigravityblue}{rgb}{0.0, 0.45, 0.74}

\newcommand{\commentary}[1]{{\small\itshape\color{commentgray} #1}}
\newcommand{\antigravity}[1]{{\small\color{antigravityblue}\textbf{Antigravity:} #1}}

% Настройка ссылок
\hypersetup{
    colorlinks=true,
    linkcolor=blue,
    filecolor=magenta,      
    urlcolor=cyan,
    pdftitle={Глубокий разбор: Метрические методы},
}

\title{Глубокий разбор: Метрические методы обучения}
\author{Твой персональный гид по ML}
\date{\today}

\begin{document}

\section{Алгоритм k-ближайших соседей (k-NN)}

\subsubsection{Математическая формулировка}
Пусть задана обучающая выборка $X^\ell = \{(x_1, y_1), \dots, (x_\ell, y_\ell)\}$, где $x_i \in \mathbb{R}^d$ — векторы признаков, а $y_i$ — целевые значения. Для произвольного объекта $x$ определим порядок близости объектов обучающей выборки:
$$\rho(x, x_{(1)}) \le \rho(x, x_{(2)}) \le \dots \le \rho(x, x_{(\ell)})$$
где $\rho(x, z)$ — функция расстояния (метрика), а $x_{(i)}$ — $i$-й ближайший сосед объекта $x$.

\commentary{Важное условие: Использование метрик требует обязательной \textbf{нормировки признаков}, так как признаки с большими абсолютными значениями будут доминировать в расчете расстояния.}

\subsubsection{Классификация и Регрессия}
\begin{itemize}
    \item \textbf{Классификация}: Ответ — наиболее часто встречающийся класс среди $k$ соседей (мода). При взвешенном голосовании:
    $a(x) = \text{argmax}_{y \in Y} \sum_{i=1}^k w_i [y_{(i)} = y]$.
    \item \textbf{Регрессия}: Ответ — среднее арифметическое ответов соседей: $a(x) = \frac{1}{k} \sum_{i=1}^k y_{(i)}$. Во взвешенном варианте: $a(x) = \frac{\sum w_i y_{(i)}}{\sum w_i}$.
\end{itemize}

\subsubsection{Выбор параметров}
\begin{itemize}
    \item \textbf{Число соседей $k$}: 
    \begin{itemize}
        \item При \textbf{малых $k$} модель имеет \textbf{низкое смещение (bias)}, но \textbf{высокий разброс (variance)}. Риск переобучения.
        \item При \textbf{больших $k$} модель имеет \textbf{высокое смещение}, но \textbf{низкий разброс}. Риск недообучения.
    \end{itemize}
    \item \textbf{Методы выбора}: Обычно $k$ подбирают на кросс-валидации или методом Leave-One-Out. Также используют «метод локтя» (анализ графика ошибки).
\end{itemize}

\commentary{Выбор $k$ — это классический поиск компромисса между недообучением и переобучением. Обычно $k$ подбирают на кросс-валидации.}

\section{Метрики и функции расстояния}

\begin{itemize}
    \item \textbf{Семейство Минковского ($L_p$)}: Общая формула: $\rho(x, z) = \left( \sum_{i=1}^d |x_i - z_i|^p \right)^{1/p}$. 
    \begin{itemize}
        \item \textbf{$L_1$ (Manhattan)}: $\sum_{i=1}^d |x_i - z_i|$. Менее чувствительна к выбросам, часто лучше работает в пространствах высокой размерности.
        \item \textbf{$L_2$ (Euclidean)}: $\sqrt{\sum_{i=1}^d (x_i - z_i)^2}$. Стандартное расстояние «по прямой». Чувствительно к большим отклонениям в отдельных признаках.
        \item \textbf{$L_\infty$ (Chebyshev)}: $\max_i |x_i - z_i|$. Расстояние определяется максимальным различием по одной из координат.
    \end{itemize}
    \item \textbf{Косинусное расстояние}: 
    $d_{\cos}(x, z) = 1 - \frac{\langle x, z \rangle}{\|x\| \|z\|} = 1 - \frac{\sum x_i z_i}{\sqrt{\sum x_i^2} \sqrt{\sum z_i^2}}$. 
    Измеряет угол между векторами, игнорируя их норму (длину). Широко применяется в анализе текстов и NLP.
    \item \textbf{Расстояние Жаккара}: 
    $d_J(A, B) = 1 - \frac{|A \cap B|}{|A \cup B|}$. 
    Используется для сравнения множеств или бинарных признаков (например, наличие слов в документе). Отражает долю общих элементов от их общего количества.
\end{itemize}

\commentary{Помните, что перед использованием $L_p$ метрик данные \textbf{обязательно} нужно масштабировать, иначе признак с большим диапазоном значений «поглотит» все остальные.}

\section{Взвешенный k-NN}

\subsection{Веса по рангу (порядковому номеру)}
Вес соседа зависит не от расстояния до него, а от его места в списке ближайших соседей.
\begin{itemize}
    \item \textbf{Линейно убывающие веса}: 
    \[ w_i = \frac{k+1-i}{k} \]
    \item \textbf{Экспоненциально убывающие веса}: 
    \[ w_i = q^i, \quad 0 < q < 1 \]
\end{itemize}

\subsection{Веса от расстояния}
Вес вычисляется непосредственно как функция от расстояния $\rho(x, x_{(i)})$.
\begin{itemize}
    \item \textbf{Обратное расстояние}: 
    \[ w_i = \frac{1}{\rho(x, x_{(i)}) + \epsilon} \]
    \item \textbf{Относительное расстояние}: 
    \[ w_i = \frac{\rho(x, x_{(k+1)}) - \rho(x, x_{(i)})}{\rho(x, x_{(k+1)}) - \rho(x, x_{(1)})} \]
\end{itemize}

\subsection{Ядерное сглаживание}
Вес соседа определяется через функцию ядра $K(z)$:
\[ w_i = K\left(\frac{\rho(x, x_i)}{h}\right) \]

\begin{itemize}
    \item \textbf{Типы ядер}:
    \begin{itemize}
        \item \textbf{Епанечникова}: \[ K(z) = \frac{3}{4} (1 - z^2) [|z| \le 1] \]
        \item \textbf{Квартическое}: \[ K(z) = \frac{15}{16} (1 - z^2)^2 [|z| \le 1] \]
        \item \textbf{Гауссовское}: \[ K(z) = \frac{1}{\sqrt{2\pi}} e^{-\frac{1}{2}z^2} \]
    \end{itemize}
\end{itemize}

\subsection{Ядерная регрессия (Формула Надарая-Ватсона)}
Это обобщение k-NN для задачи регрессии, где предсказание — это взвешенное среднее ответов всех объектов обучающей выборки:
\[ a(x) = \frac{\sum_{i=1}^n y_i K\left(\frac{\rho(x, x_i)}{h}\right)}{\sum_{i=1}^n K\left(\frac{\rho(x, x_i)}{h}\right)} \]

\subsection{Оптимизация поиска ближайшего соседа}
Для работы с большими данными ($N \gg 1$) используются методы, позволяющие избежать полного перебора объектов.

\begin{itemize}
    \item \textbf{Точные методы}:
    \begin{itemize}
        \item \textbf{k-d деревья}: Разбивают пространство по признакам. Эффективны только при малых $d$ ($< 20$). При высокой размерности деградируют до $O(N)$.
        \item \textbf{Ball trees}: Разбивают пространство на гиперсферы. Работают стабильнее k-d деревьев в пространствах средней размерности.
    \end{itemize}
    \item \textbf{Приближенные методы (ANN)}:
    \begin{itemize}
        \item \textbf{LSH (Locality-Sensitive Hashing)}: Хеширование, при котором близкие объекты с высокой вероятностью попадают в один бакет. Поиск идет только внутри бакета.
        \item \textbf{HNSW (Hierarchical Navigable Small World)}: Построение графа, где вершины — объекты, а ребра соединяют соседей. Иерархическая структура позволяет быстро выходить в нужную область пространства и уточнять результат.
        \item \textbf{Inverted File Index (IVF)}: Пространство делится на кластеры (ячейки Вороного), и поиск ведется только в ближайших к запросу кластерах.
    \end{itemize}
\end{itemize}

\antigravity{В реальном продакшене (Faiss, HNSW) почти всегда используются именно ANN, так как точный поиск на миллионах объектов невозможен.}

\end{document}
